% --- SECTION IV: PROPOSED FRAMEWORK (Mobility-Aware Predictive Placement) ---
\section{Mobility-Aware Predictive Placement}
\label{sec:method}

This section presents the proposed framework that combines: (i) an \emph{offline} lower bound computed from a linear relaxation of the assignment ILP, (ii) a deep mobility predictor that outputs a \emph{probability map} over candidate edge nodes, and (iii) an \emph{online} scheduling heuristic that performs mobility-aware placement while respecting resource constraints.

\subsection{Offline Bound via Linear Relaxation}
\label{subsec:lp_bound}

At each decision step \(t\), the assignment problem in Section~III-C is an ILP with binary variables \(x_{u,n}(t)\in\{0,1\}\).
To obtain a tractable lower bound for benchmarking online policies, we relax integrality:
\[
x_{u,n}(t)\in[0,1],\qquad \forall u\in\mathcal{U}(t),\ \forall n\in\mathcal{N}.
\]
All other constraints (one-to-one assignment and capacity constraints) are kept unchanged. The resulting linear program (LP) yields an objective value that is a \emph{lower bound} on the optimal ILP objective (since it optimizes over a superset of feasible solutions).

In our implementation, the relaxed assignment matrix is solved with a convex optimization solver. For numerical stability, the turnaround-time matrix can be normalized by its maximum entry without changing the argmin structure. After solving the relaxation, a discrete assignment is derived by selecting, for each user \(u\), the node with the largest relaxed weight and then re-checking feasibility; infeasible selections fall back to the central node. This produces an efficient approximation while the LP objective remains a useful offline bound for comparison.

\subsection{Deep Mobility Predictor}
\label{subsec:predictor}

The framework uses a GRU-based trajectory predictor trained on T-Drive/TaxiD trajectories. We keep the architectural discussion brief and emphasize the interface: \emph{inputs}, \emph{outputs}, and the \emph{probability map} used by the scheduler.

\paragraph{Input (feature sequence).}
For each user \(u\), the central node maintains a rolling history of the most recent \(L\) observations (default \(L=20\)). Each history record is represented in metric coordinates (meters) and includes a timestamp and a feature vector.
Depending on availability, the history can contain:
\begin{itemize}
  \item \textbf{Core kinematic/temporal features} computed during offline preprocessing: position \((x,y)\), speed \(v\), acceleration \(a\), \(\Delta v\), heading change \(\Delta\theta\), and time encodings (e.g., \(\sin/\cos\) of time-of-day and day-of-week), plus stop/dwell indicators.
  \item \textbf{Optional graph-context features} (Phase B): node degree and junction flag, when map-matched road context is available.
\end{itemize}
At inference time, missing optional features are filled with zeros and the full feature vector is standardized using the scaler saved with the model artifact.

\paragraph{Output (future waypoints and horizons).}
The deployed model runs in a \texttt{curv\_step} mode: given the lookback window, it predicts a sequence of stepwise curvilinear displacements \(\{\Delta s_k\}_{k=1}^{K}\) for \(K=10\) future minutes. These displacements are then rolled out into absolute future positions \(\hat{\mathbf{p}}_u(h)=(\hat{x}_u(h),\hat{y}_u(h))\) for a set of horizons \(h\in\{1,3,5,10\}\) minutes. When a road graph is configured, rollout can be constrained along the road network; otherwise, a heading-based free-space rollout is used.

\paragraph{Probability map (user-to-cloudlet likelihoods).}
Let \(\mathbf{p}_e\) be the (fixed) position of edge node \(e\in\mathcal{E}\) in meters. For each horizon \(h\), the predictor converts the predicted waypoint \(\hat{\mathbf{p}}_u(h)\) to a probability distribution over edge nodes by applying a softmax over negative Euclidean distances:
\[
P_{u,e}(h)=\frac{\exp\left(-\frac{\|\hat{\mathbf{p}}_u(h)-\mathbf{p}_e\|_2}{T}\right)}
{\sum_{e'\in\mathcal{E}}\exp\left(-\frac{\|\hat{\mathbf{p}}_u(h)-\mathbf{p}_{e'}\|_2}{T}\right)},
\]
where \(T\) is a temperature parameter (default \(T=50\)) controlling the sharpness of the distribution. Optionally, a maximum radius can be used to downweight faraway cloudlets; if all cloudlets fall outside the radius, the distribution reverts to an unmasked softmax. The final output is a \emph{probability tensor} \(P_u\in[0,1]^{|\mathcal{E}|\times |H|}\), where \(H=\{1,3,5,10\}\).

For scalability with many users, batch inference is supported: the predictor performs a single forward pass for all eligible users in \(\mathcal{U}(t)\), then computes their probability maps.

\subsection{Online Predictive Scheduler}
\label{subsec:online_scheduler}

The online scheduler is a lightweight heuristic that uses the probability map to guide placement decisions while enforcing resource feasibility. It operates in two closely related modes: \emph{immediate predictive assignment} and \emph{prewarm-only planning}.

\paragraph{Step 1: Obtain probability map.}
At step \(t\), the scheduler partitions users into:
(i) users with sufficient history length (\(\geq L\)) to run prediction, and
(ii) users without sufficient history, which fall back to greedy distance-based assignment.
For eligible users, the scheduler queries the predictor to obtain \(P_u\).

\paragraph{Step 2: Expected cold-start risk (execution-aware signal).}
Cold starts depend on the relationship between the chosen node and the previously executed node for the same logical function/user. In the simulator, this signal is tracked as the last executed node ID, and cold/warm execution time is assigned accordingly (in simulated execution mode).
Given \(P_u\) at a chosen horizon \(h^\star\), an execution-aware proxy for cold-start risk can be expressed as:
\[
R_u(t) = 1 - P_{u,e_{\text{prev}}}(h^\star),
\]
where \(e_{\text{prev}}\) is the last executed node for user \(u\). In the current implementation, the scheduler does not explicitly optimize \(R_u(t)\) as a separate term; instead, cold/warm effects are reflected by the execution layer during evaluation, and by the prewarm-only mode described below.

\paragraph{Step 3: Probabilistically greedy decision rule.}
For each user \(u\) with a valid probability map, the scheduler selects a target horizon \(h^\star\) (default 5 minutes) and sorts edge nodes by descending probability \(P_{u,e}(h^\star)\). It then chooses the first edge node that satisfies resource constraints (CPU/memory health thresholds and available memory capacity). If no edge node is feasible, the user is assigned to the central node. Users without a probability map (e.g., insufficient history or inference failure) are assigned by greedy nearest-feasible selection.

\paragraph{Pre-warming logic (planning ahead).}
To model pre-warming under slow wall-clock simulation steps (where Docker warm TTL is not meaningful), the scheduler supports a \emph{prewarm-only} mode:
\begin{enumerate}
  \item The scheduler keeps serving requests on the current assignment.
  \item On periodic planning ticks (every \(I\) steps), it computes a planned node \(e^{\text{plan}}_u\) from the probability map and stores a planned switch step \(t^{\text{plan}} = t + \Delta\) (lead time \(\Delta\) steps).
  \item When \(t=t^{\text{plan}}\), the scheduler applies the switch to \(e^{\text{plan}}_u\). On that switch step, the execution layer treats the request as \emph{prewarmed} (warm start) when simulating cold/warm behavior.
\end{enumerate}
This design captures the intended effect of ``starting containers early'' using an analytical warm/cold model, while avoiding extra network/Docker overhead in large-scale experiments.

